\documentclass[11pt]{article}
\usepackage[margin=1in]{geometry}
\usepackage{amsmath,amssymb,amsthm}
\usepackage{hyperref}
\newtheorem{theorem}{Theorem}
\newtheorem{lemma}[theorem]{Lemma}
\newtheorem{proposition}[theorem]{Proposition}
\newtheorem{corollary}[theorem]{Corollary}
\theoremstyle{definition}
\newtheorem{definition}[theorem]{Definition}
\theoremstyle{remark}
\newtheorem{remark}[theorem]{Remark}
\newcommand{\Q}{\mathbb{Q}}
\newcommand{\Z}{\mathbb{Z}}
\newcommand{\C}{\mathbb{C}}
\newcommand{\OO}{\mathcal{O}}
\newcommand{\disc}{\operatorname{disc}}
\newcommand{\Tr}{\operatorname{Tr}}
\newcommand{\Nm}{\operatorname{N}}
\newcommand{\Cl}{\operatorname{Cl}}

\title{Disproof of Conjecture 00000004310:\\
imaginary quadratic fields of the same discriminant\\ have isomorphic unit groups}
\author{jilint777}
\date{2026-10-04}

\begin{document}
\sloppy
\maketitle

\begin{abstract}
Conjecture 00000004310 asserts that there exist two imaginary quadratic fields of the same
discriminant whose class groups have isomorphic $3$-parts while their unit groups are not
isomorphic, and that the smallest such discriminant has seven digits. This is false, whatever
the class-group condition means. The discriminant determines an imaginary quadratic field: every
such field is $\Q(\sqrt m)$ for a unique squarefree $m<0$, its discriminant is $m$ or $4m$, and
$m\mapsto\disc m$ is injective. So two fields with the same discriminant are equal, and their
unit groups are equal. A second, independent argument: the unit group of an imaginary quadratic
field, and more generally of an imaginary quadratic order, is cyclic of order $w(D)$, where
$w(-4)=4$, $w(-3)=6$ and $w(D)=2$ otherwise, so it is a function of the discriminant $D$ alone.
There is therefore no witness discriminant at all, and in particular no smallest one with seven
digits. The Lean~4 development (core library only) defines squarefree integers, the
discriminant, the rings of integers $\OO_m$ and the quadratic orders of every discriminant as
rings of pairs, their unit groups and isomorphisms of unit groups from scratch. It proves
injectivity of the discriminant, the complete list of units, and the negation of the claim, for
an arbitrary class-group predicate. A standard-library Python script checks the same facts
numerically by other methods, including the reading with binary quadratic forms.
\end{abstract}

\section{The statement}
The conjecture reads:
\begin{quote}
Definition: The $p$-part of the class group and the unit group are two invariants of fields of
the same discriminant. Conjecture: There exist two imaginary quadratic fields of the same
discriminant whose $3$-parts of the class groups are isomorphic while their unit groups are not,
and the smallest discriminant is a seven-digit value.
\end{quote}
The Chinese version says the same: there exist two imaginary quadratic fields with the same
discriminant whose class groups have isomorphic $3$-parts and whose unit groups are not
isomorphic, and the minimal discriminant is a seven-digit number.

The conjecture is a conjunction of two clauses:
\begin{enumerate}
\item[(E)] there are imaginary quadratic fields $K_1,K_2$ with $D_{K_1}=D_{K_2}$,
$\Cl(K_1)[3^\infty]\cong\Cl(K_2)[3^\infty]$ and $\OO_{K_1}^\times\not\cong\OO_{K_2}^\times$;
\item[(M)] the smallest $|D|$ for which (E) holds has seven digits.
\end{enumerate}
We show that (E) is false. Then (M) is false as well, since there is no witness discriminant at
all. Refuting (E) suffices to refute the conjecture.

\section{Background}
We use the following standard facts; see D.~A.~Marcus, \emph{Number Fields}, 2nd ed., Springer
2018, Chapter~2 (quadratic fields, their rings of integers and discriminants), and D.~A.~Cox,
\emph{Primes of the Form $x^2+ny^2$}, 2nd ed., Wiley 2013, \S2 (forms) and \S7 (orders).

\begin{enumerate}
\item[(F1)] \emph{Classification.} Every quadratic field is $\Q(\sqrt m)$ for a unique squarefree
integer $m\notin\{0,1\}$. It is imaginary iff $m<0$. Uniqueness: $\Q(\sqrt m)=\Q(\sqrt{m'})$
forces $m/m'$ to be a square in $\Q$, hence $m=m'$ for squarefree $m,m'$. This holds for
subfields of $\C$ and for abstract fields up to isomorphism.
\item[(F2)] \emph{Ring of integers.} For squarefree $m$, $\OO_K=\Z[\sqrt m]$ if
$m\equiv 2,3\pmod 4$ and $\OO_K=\Z[\tfrac{1+\sqrt m}{2}]$ if $m\equiv1\pmod 4$.
\item[(F3)] \emph{Discriminant.} $D_K=4m$ if $m\equiv2,3\pmod4$ and $D_K=m$ if $m\equiv1\pmod4$.
\item[(F4)] \emph{Orders.} For every $D<0$ with $D\equiv0,1\pmod 4$ there is exactly one
imaginary quadratic order of discriminant $D$ up to isomorphism, namely
$\OO_D=\Z\big[\tfrac{D+\sqrt D}{2}\big]$. Every order in an imaginary quadratic field is of this
form (Cox, \S7.A: an order of conductor $f$ is $\Z+f\OO_K$, of discriminant $f^2D_K$).
\end{enumerate}

\section{The discriminant determines the field}
For an integer $m$ put $\disc m=m$ if $m\equiv1\pmod 4$ and $\disc m=4m$ otherwise.

\begin{lemma}\label{lem:inj}
The map $m\mapsto\disc m$ is injective on $\Z$. On squarefree $m$ its values satisfy
$\disc m\equiv1\pmod 4$ or $\disc m\equiv 8,12\pmod{16}$, and $m$ is recovered as
$m=\disc m$ if $\disc m\equiv1\pmod4$ and $m=\disc m/4$ otherwise.
\end{lemma}
\begin{proof}
Suppose $\disc m_1=\disc m_2$. If $m_1\equiv m_2\equiv1\pmod4$, then $m_1=m_2$ directly. If
neither is $\equiv1$, then $4m_1=4m_2$. If $m_1\equiv1$ and $m_2\not\equiv1$, then
$m_1=4m_2\equiv0\pmod4$, a contradiction; the remaining case is symmetric. For squarefree $m$ we
have $4\nmid m$. So if $m\not\equiv1\pmod 4$ then $m\equiv2,3\pmod4$ and
$4m\equiv8,12\pmod{16}$. The recovery formula follows because $4m\equiv0\not\equiv1\pmod4$.
\end{proof}

\begin{theorem}\label{thm:same-field}
If $K_1,K_2$ are imaginary quadratic fields with $D_{K_1}=D_{K_2}$, then $K_1\cong K_2$. If they
are subfields of $\C$, then $K_1=K_2$. In either case
$\OO_{K_1}^\times\cong\OO_{K_2}^\times$.
\end{theorem}
\begin{proof}
By (F1) write $K_i\cong\Q(\sqrt{m_i})$ with $m_i<0$ squarefree. By (F3),
$\disc m_1=D_{K_1}=D_{K_2}=\disc m_2$, so $m_1=m_2$ by Lemma~\ref{lem:inj}. Hence
$K_1\cong\Q(\sqrt{m_1})\cong K_2$. Inside $\C$, an imaginary quadratic field is determined by the
square class of $m$, so $K_1=K_2$. A field isomorphism maps the ring of integers onto the ring of
integers, so it restricts to an isomorphism of unit groups.
\end{proof}

\section{The unit groups depend only on the discriminant}
This section gives a second argument that does not use Lemma~\ref{lem:inj}, and it also covers
orders.

\paragraph{Quadratic rings of pairs.}
For integers $t,n$ let $R_{t,n}=\Z[\omega]$, $\omega^2=t\omega-n$. Its elements are
$a+b\omega$, $a,b\in\Z$, with
\[
(a+b\omega)(c+d\omega)=(ac-nbd)+(ad+bc+tbd)\,\omega .
\]
The ring has the conjugation $\overline{a+b\omega}=(a+tb)-b\omega$, the trace
$\Tr(a+b\omega)=2a+tb$, and the norm $\Nm(a+b\omega)=(a+b\omega)\overline{(a+b\omega)}=a^2+tab+nb^2$.
The norm is multiplicative and satisfies
\begin{equation}\label{eq:4N}
4\Nm(a+b\omega)=(2a+tb)^2+(4n-t^2)\,b^2 .
\end{equation}
The discriminant of the basis $(1,\omega)$ is
$\det\big(\Tr(e_ie_j)\big)=2(t^2-2n)-t^2=t^2-4n$.

\paragraph{Orders and maximal orders.}
For $D\equiv0,1\pmod4$ put $t=D\bmod 2\in\{0,1\}$ and $n=(t-D)/4$. Then $t^2-4n=D$, and
$\omega=(t+\sqrt D)/2$ satisfies $\omega^2=t\omega-n$, so $R_{t,n}\cong\OO_D$ is the order of
discriminant $D$ in (F4). For squarefree $m$, the parameters of (F2) are $(t,n)=(1,(1-m)/4)$ when
$m\equiv1\pmod4$, with $\omega=(1+\sqrt m)/2$, and $(t,n)=(0,-m)$ otherwise, with $\omega=\sqrt m$.
These are exactly the parameters attached to $D=\disc m$, so $\OO_{\Q(\sqrt m)}=\OO_{\disc m}$.

The map $a+b\omega\mapsto(x,y)$, where $x+y\sqrt m=2(a+b\omega)$, identifies $R_{t,n}$ with the
elements $(x+y\sqrt m)/2$ of $\Q(\sqrt m)$ and respects multiplication. Formula (F2) is the
following statement.

\begin{lemma}[integrality]\label{lem:integral}
Let $m\not\equiv0\pmod4$ and $x,y\in\Z$. The element $(x+y\sqrt m)/2$, which has trace
$x\in\Z$, has norm $(x^2-my^2)/4\in\Z$ iff it lies in $\OO_m$; that is, iff $x\equiv y\pmod2$
when $m\equiv1\pmod4$, and iff $x,y$ are both even otherwise.
\end{lemma}
\begin{proof}
A square $z^2$ is $\equiv0\pmod4$ if $z$ is even and $\equiv1\pmod4$ if $z$ is odd, so $x^2-my^2\equiv[x\text{ odd}]-m[y\text{ odd}]\pmod 4$.
If $y$ is even this is $\equiv0$ iff $x$ is even. If $y$ is odd it is $\equiv 1-m$ or $-m$,
according to the parity of $x$. Since $4\nmid m$, the value is $0$ exactly when $x$ is odd and
$m\equiv1\pmod4$.
\end{proof}
Every algebraic integer $\alpha=u+v\sqrt m\in\Q(\sqrt m)$ has this form: $2u=\Tr\alpha\in\Z$, and
$(2u)^2-m(2v)^2=4\Nm\alpha\in\Z$, so $m(2v)^2\in\Z$, and $2v\in\Z$ because $m$ is squarefree.
Together with Lemma~\ref{lem:integral} this is the proof of (F2).

\begin{theorem}[units]\label{thm:units}
Let $D<0$, $D\equiv0,1\pmod4$, and let $(t,n)$ be as above. The units of $\OO_D=R_{t,n}$ are the
elements of norm $1$. They are
\begin{itemize}
\item $\pm1,\pm\omega$ for $D=-4$ (here $\omega=i$),
\item $\pm1,\pm\omega,\pm(\omega-1)$ for $D=-3$ (here $\omega=(1+\sqrt{-3})/2$, $\omega^2=\omega-1$),
\item $\pm1$ for every other $D$.
\end{itemize}
In each case $\OO_D^\times$ is cyclic of order $w(D)$, where $w(-4)=4$, $w(-3)=6$ and $w(D)=2$
otherwise. It is generated by $\omega$, by $\omega$, and by $-1$, respectively.
\end{theorem}
\begin{proof}
Since $4n-t^2=-D>0$, \eqref{eq:4N} shows that $\Nm\ge0$. If $xy=1$, then
$\Nm(x)\Nm(y)=1$ forces $\Nm(x)=1$. Conversely, if $\Nm x=1$ then $x\bar x=1$. Now let
$\Nm(a+b\omega)=1$. By \eqref{eq:4N}, $4=(2a+tb)^2+|D|b^2$ with $|D|\ge3$, so $|b|\le1$ and
$|2a+tb|\le2$, hence $|a|\le1$. If $|D|>4$ then $b=0$ and $a=\pm1$. For $D=-4$ the equation is
$a^2+b^2=1$, and for $D=-3$ it is $a^2+ab+b^2=1$. Enumerating the nine pairs $(a,b)$ gives the
lists. The powers of $\omega$ for $D=-3$ are
$1,\omega,\omega-1,-1,-\omega,1-\omega$, and $\omega^6=1$.
\end{proof}

\begin{corollary}\label{cor:w}
For imaginary quadratic fields, $\OO_{\Q(\sqrt m)}^\times\cong\Z/w(\disc m)$. If $w(D_1)=w(D_2)$,
then the unit groups of the orders of discriminants $D_1,D_2$ are isomorphic. In particular,
equal discriminants give isomorphic unit groups. For orders of a fixed discriminant this holds
without using Lemma~\ref{lem:inj}.
\end{corollary}

\begin{theorem}\label{thm:main}
Clause (E) is false. More precisely, for every predicate $C(K_1,K_2)$ standing in for
``$\Cl(K_1)[3^\infty]\cong\Cl(K_2)[3^\infty]$'', there are no imaginary quadratic fields with
$D_{K_1}=D_{K_2}$, $C(K_1,K_2)$ and $\OO_{K_1}^\times\not\cong\OO_{K_2}^\times$. Consequently
there is no witness discriminant, and clause (M) is false.
\end{theorem}
\begin{proof}
This follows from Theorem~\ref{thm:same-field}, or independently from Corollary~\ref{cor:w}.
\end{proof}

\section{Readings of the statement}\label{sec:readings}
\begin{enumerate}
\item \textbf{``Two fields'' as subfields of $\C$ or as fields up to isomorphism.} In both cases
Theorem~\ref{thm:same-field} gives $K_1\cong K_2$, hence isomorphic unit groups.
\item \textbf{``Two'' fields that are required to be distinct.} Then there is not even a pair
of distinct imaginary quadratic fields with the same discriminant, so (E) is false even before
units are considered. If $K_1=K_2$ is allowed, the unit groups are equal.
\item \textbf{The class-group condition.} It is one conjunct of an existential, and we refute
the existential for an arbitrary predicate in its place. Whatever ``$3$-part'' means (Sylow
$3$-subgroup, $3$-rank, $\Cl/\Cl^3$), the argument is unaffected. When $K_1=K_2$ it holds
trivially.
\item \textbf{Orders and binary quadratic forms.} If ``fields of discriminant $D$'' is read as
imaginary quadratic orders of discriminant $D$, then by (F4) there is only one, and by
Theorem~\ref{thm:units} its unit group is $\Z/w(D)$, a function of $D$. If it is read as
primitive positive definite binary quadratic forms of discriminant $D$, the form class group
$C(D)\cong\operatorname{Pic}(\OO_D)$ is a function of $D$. The ``unit group'' of a form $f$, its
group of proper automorphs in $\mathrm{SL}_2(\Z)$, is in bijection with the solutions of
$t^2-Du^2=4$ via $(t,u)\mapsto\begin{pmatrix}(t-bu)/2&-cu\\au&(t+bu)/2\end{pmatrix}$ (the classical
description of automorphs; see Cox, \S7, and D.~A.~Buell, \emph{Binary Quadratic Forms},
Springer 1989, Ch.~3). So it is cyclic of
order $w(D)$ for every form of discriminant $D$. Again nothing depends on more than $D$.
\item \textbf{Absolute value of the discriminant.} Both fields are imaginary, so their
discriminants are negative, and $|D_1|=|D_2|$ gives $D_1=D_2$.
\item \textbf{Dropping ``of the same discriminant''.} Then pairs with isomorphic $3$-parts and
non-isomorphic unit groups do exist, but they are tiny. For example $\Q(\sqrt{-3})$,
$\Q(i)$ and $\Q(\sqrt{-2})$ all have class number $1$ (trivial $3$-parts), with unit groups of
orders $6,4,2$. The discriminants involved are $-3,-4,-8$, with one digit, so clause (M), the
seven-digit minimum, fails under this reading too. Moreover, every witness pair must contain
$\Q(i)$ or $\Q(\sqrt{-3})$, since all other unit groups are $\{\pm1\}$.
\item \textbf{Real quadratic fields of the same discriminant.} The word ``imaginary'' is
explicit, but the claim fails for real fields too. Every real quadratic field is $\Q(\sqrt m)$
for a unique squarefree $m>1$, its discriminant is again $m$ or $4m$ by (F3), and
Lemma~\ref{lem:inj} holds on all of $\Z$. So two real quadratic fields with the same
discriminant are equal. In any case, by Dirichlet's unit theorem every real quadratic unit group
is $\{\pm1\}\times\varepsilon^{\Z}\cong\Z/2\times\Z$, so all of them are isomorphic.
\item \textbf{A real and an imaginary field with the same $|D|$.} Here the unit groups always
differ ($\Z/2\times\Z$ against a finite group), so such pairs can be witnesses. They are
small. Take $\Q(\sqrt2)$ and $\Q(\sqrt{-2})$, with $D=8$ and $D=-8$. Both have class number
$1$: the Minkowski bounds are $\tfrac12\sqrt8\approx1.41$ and
$\tfrac2\pi\sqrt8\approx1.80$, both below $2$, so every ideal class contains $\OO_K$ and the
$3$-parts are trivial. The unit groups are $\Z/2\times\Z$ and $\Z/2$. Here $|D|=8$ has one
digit, so clause (M) fails under this reading too.
\end{enumerate}

\begin{remark}[on ``seven digits'']
The seven-digit figure possibly echoes the classical computations of imaginary quadratic fields
whose class group has large $3$-rank (Diaz y Diaz, 1974). There the smallest discriminants with
$3$-rank $3$ are of seven-digit size. That is a statement about the $3$-part of the class group
alone. It cannot be rescued by keeping the unit-group condition: a separation of unit groups
between imaginary quadratic fields forces $\Q(i)$ or $\Q(\sqrt{-3})$ (discriminant $-4$ or
$-3$) into the pair, by Theorem~\ref{thm:units}. A separation between real fields is
impossible, and a real/imaginary pair already occurs at $|D|=8$. So no repair of the statement
that keeps the unit-group condition has a seven-digit minimum.
\end{remark}

\section{The Lean formalization}
The directory \texttt{lean4/} contains a Lean~4.19.0 project that uses the core library only.
All objects are defined from scratch.
\begin{itemize}
\item \texttt{Squarefree m} ($d^2\mid m\Rightarrow d=\pm1$), \texttt{IsImagQuadParam m}
($m<0$ and squarefree), and \texttt{disc m}. \texttt{SqfCheck} is a decidable test proved sound
(\texttt{squarefree\_of\_check}). It gives \texttt{param\_neg1}, \texttt{param\_neg2},
\texttt{param\_neg3}, \dots; \texttt{not\_param\_neg4}.
\item \texttt{disc\_injective}, \texttt{disc\_fundamental} (values $\equiv1\bmod4$ or
$\equiv8,12\bmod16$), and \texttt{mOfDisc\_disc} (explicit inverse): Lemma~\ref{lem:inj}.
\item \texttt{QR} and \texttt{mul t n}: the ring $R_{t,n}$ of pairs. Proven:
\texttt{mul\_comm}, \texttt{mul\_assoc}, \texttt{one\_mul}, \texttt{mul\_one},
\texttt{mul\_add}, \texttt{omega\_sq} ($\omega^2=t\omega-n$), \texttt{norm\_mul},
\texttt{mul\_conj}, \texttt{four\_norm} (equation~\eqref{eq:4N}), and \texttt{discBasis\_eq}
($\det\Tr(e_ie_j)=t^2-4n$).
\item \texttt{fT m}, \texttt{fN m}: the parameters of $\OO_m$ in (F2). \texttt{sqrtM\_sq}:
$\sqrt m\in\OO_m$. \texttt{field\_discBasis}: the discriminant of $\OO_m$ is \texttt{disc m}.
\texttt{halfCoords\_mul}: under $a+b\omega\mapsto(x+y\sqrt m)/2$ the multiplication is that of
$\Q(\sqrt m)$. \texttt{integral\_iff}: Lemma~\ref{lem:integral}.
\item \texttt{tD D}, \texttt{nD D}: the order of discriminant $D$, with
\texttt{discBasis\_tD}. \texttt{fT\_eq}, \texttt{fN\_eq}: $\OO_m$ is the order of discriminant
\texttt{disc m}.
\item \texttt{IsUnitQ} (has an inverse), \texttt{UnitGrp} (the subtype with its
multiplication), and \texttt{isUnit\_iff\_norm}.
\texttt{units\_classification}: for every valid $D<0$, $x$ is a unit iff
$x\in$ \texttt{unitList D}. \texttt{unitList\_length} ($=w(D)$), \texttt{unitList\_nodup}, and
\texttt{unitList\_eq\_pows}: the list is $1,g,\dots,g^{w-1}$ with $g^{w}=1$. Together these say
that the unit group is cyclic of order $w(D)$ (Theorem~\ref{thm:units}).
\item \texttt{field\_units\_pm\_one}, \texttt{field\_units\_neg1}, \texttt{field\_units\_neg3},
and \texttt{field\_units\_card}: the field case.
\item \texttt{MulIso}: bijections with inverse that preserve multiplication.
\texttt{orderUnits\_iso\_of\_w\_eq} and \texttt{fieldUnits\_iso\_of\_disc\_eq}:
Corollary~\ref{cor:w}. The latter does not use \texttt{disc\_injective}.
\item \texttt{Claim C3}: $\exists\,m_1,m_2$, \texttt{IsImagQuadParam} $m_1$,
\texttt{IsImagQuadParam} $m_2$, \texttt{disc} $m_1=$ \texttt{disc} $m_2$, \texttt{C3} $m_1\,m_2$,
and $\neg$\,\texttt{Nonempty (MulIso (FieldUnits} $m_1$\texttt{) (FieldUnits} $m_2$\texttt{))}.
Here \texttt{C3} is an arbitrary predicate standing for the class-group condition.
\texttt{conjecture\_00000004310\_false : }$\neg$\,\texttt{Claim C3} (via injectivity), and
\texttt{conjecture\_00000004310\_false'} (via $w$). \texttt{no\_witness\_discriminant} and
\texttt{smallest\_seven\_digit\_false} refute clause (M). \texttt{order\_claim\_false} covers
the reading with orders.
\item Non-vacuity: \texttt{fieldUnits\_not\_iso} proves
$\OO_{-1}^\times\not\cong\OO_{-2}^\times$ by a pigeonhole argument, so the condition ``not
isomorphic'' is not empty. \texttt{weakened\_reading\_small\_witness} records the pair
$(-1,-2)$ with discriminants $-4,-8$ for reading~6. \texttt{hypotheses\_satisfiable} and
\texttt{param\_*} show that the field hypotheses are satisfiable.
\end{itemize}
Polynomial identities are proved by distributing, sorting monomials with
\texttt{simp only [Int.mul\_comm, \dots]}, and closing with \texttt{omega} (the macro
\texttt{poly\_omega}). The build takes about ten seconds. The \texttt{\#print axioms} report
shows only \texttt{propext}, \texttt{Quot.sound} and \texttt{Classical.choice}. There is no
\texttt{sorry} and no \texttt{native\_decide}.

\paragraph{What is not in Lean.} Lean does not construct $\Q(\sqrt m)$ or the ring of all
algebraic integers. Facts (F1)--(F4) are the bridge from ``imaginary quadratic field'' to
``squarefree $m<0$ together with the ring \texttt{QR} with parameters \texttt{fT m},
\texttt{fN m}''. Lean proves the integer part of (F2) (\texttt{integral\_iff},
\texttt{halfCoords\_mul}) and the discriminant formula (F3) as the determinant of the trace
form. The rational step $2v\in\Z$ and the uniqueness in (F1) are argued in this report. Class
groups are not formalized; they enter only as an arbitrary predicate. Once fields are parametrised by $m$, the first Lean proof
\texttt{conjecture\_00000004310\_false} is immediate: it uses only \texttt{disc\_injective}.
The mathematical content of the disproof is the cited classification (F1)--(F4) together with
the $w(D)$ argument, which is formalized in \texttt{units\_classification},
\texttt{unitList\_eq\_pows}, \texttt{orderUnits\_iso\_of\_w\_eq} and the second proof
\texttt{conjecture\_00000004310\_false'}. The real-quadratic readings 7 and 8 are argued in
this report and checked in \texttt{verify.py}, not in Lean.

\section{Independent verification}
\texttt{verify.py} (standard library only) checks:
\begin{itemize}
\item that \texttt{disc} is injective on the $12160$ squarefree $m$ with $-20000<m<0$; that its
values are fundamental discriminants (recognised by an independent test); and that all $6079$
negative fundamental discriminants with $|D|<20000$ arise;
\item the unit counts of $\OO_m$ by brute force on a box of side 7, giving $4$, $6$ and $2$, so
the count is a function of $D$;
\item for all $4999$ discriminants $-10^4<D<0$, $D\equiv0,1\pmod4$, that $t^2-Du^2=4$ and
$\Nm(a+b\omega)=1$ each have $w(D)$ solutions, and that the unit group is cyclic;
\item for binary quadratic forms, that every reduced primitive form of discriminant $D$
($143695$ forms for $|D|<10^4$) has exactly $w(D)$ proper automorphs. The automorphs are found
through $(t,u)$ for all $D$, and by brute force over $\mathrm{SL}_2$ matrices with entries in
$[-2,2]$ for $|D|<400$;
\item a few class numbers, and that no discriminant carries two fields with different unit groups;
\item reading~6: $h(-3)=h(-4)=h(-8)=1$, giving one-digit witnesses;
\item readings~7--8: injectivity of \texttt{disc} on the squarefree $1<m<20000$; for every
squarefree $1<m<2000$, a unit of infinite order (a solution of $x^2-my^2=\pm1$, $y>0$, found by
continued fractions); and for $\Q(\sqrt{\pm2})$, $|D|=8$, Minkowski bounds below $2$.
\end{itemize}
All checks pass (see \texttt{verification.txt}).

\section{Conclusion}
An imaginary quadratic field is determined by its discriminant, and its unit group is cyclic of
order $w(D)\in\{2,4,6\}$, a function of the discriminant alone. So no two imaginary quadratic
fields of the same discriminant have non-isomorphic unit groups, whatever their class groups
are. The same holds for imaginary quadratic orders and for binary quadratic forms. The witness
set is empty, so there is no ``smallest seven-digit discriminant''. Conjecture 00000004310 is
false.

\end{document}
